\documentclass[11pt,a4paper]{article}
\usepackage[margin=2.5cm]{geometry}
\usepackage[hidelinks]{hyperref}
\usepackage{graphicx}
\usepackage{booktabs}
\usepackage{amsmath}
\usepackage{enumitem}
\usepackage[T1]{fontenc}
\usepackage[utf8]{inputenc}
\usepackage{csquotes}
\usepackage{microtype}
\setlist[itemize]{noitemsep,topsep=2pt}
\setlist[enumerate]{noitemsep,topsep=2pt}

\title{Trust Analysis of Traffic Sign Classifiers under Uncertainties\\[2mm]
Bachelor Thesis Proposal}
\author{Youssef Hesham}
\date{\today}

\begin{document}
\maketitle

\section{Motivation}
\subsection{Research field}
Perception for intelligent vehicles combines real-time detection, tracking, and classification to interpret traffic scenes and support Advanced Driver Assistance Systems and autonomous driving. Traffic Sign Recognition (TSR) is a core capability: the system must localize signs in full-frame images and then classify them into semantic categories whose meanings carry legal and safety consequences.

In real traffic, perception operates under uncertainty. Inputs vary with illumination, weather, motion blur, sensor noise, compression artifacts, lens contamination, partial occlusions, domain shift across regions and cameras, and imperfect detections that crop a sign too tightly or too loosely. Even small changes can hide salient digits, distort shapes and colors, or reduce resolution, pushing the model outside the conditions it learned from.

Uncertainty in TSR is commonly divided into aleatoric and epistemic components. Aleatoric uncertainty stems from the data itself—missing pixels, blur, glare, shadows, or occluding objects—and cannot be eliminated by training alone. Epistemic uncertainty reflects a model’s limited knowledge due to finite or biased training data; it is higher on unfamiliar styles of signs, unusual viewpoints, or rare classes, and can be reduced with better data and modeling.

Modern TSR systems use convolutional backbones such as ResNet variants for classification and efficient one‑stage detectors for proposals, tuned to meet real-time constraints on embedded hardware. While these models can reach high top‑1 accuracy on curated datasets like GTSRB, their predicted probabilities often fail to reflect true likelihoods of correctness once uncertainty increases. Overconfident errors are especially problematic in safety‑critical contexts.

Trustworthy TSR therefore requires more than accuracy. It demands calibrated confidence—probabilities that match empirical correctness—and mechanisms to act on uncertainty. Common diagnostics include reliability diagrams, Expected Calibration Error, Brier score, and Negative Log‑Likelihood. Beyond diagnostics, systems can use uncertainty to improve decisions: cross‑checking with maps, slowing down, requesting human confirmation, or abstaining when confidence is low.

Research on trustworthy perception investigates how to estimate uncertainty at test time, how to calibrate probabilities, and how to design decision policies that trade coverage for risk under strict latency budgets. Lightweight post‑hoc calibration methods such as temperature scaling or isotonic regression change only the final scores and add almost no runtime. Other approaches—entropy or energy‑based scoring, Monte‑Carlo dropout, or small deep ensembles—provide richer uncertainty signals, albeit with increasing computational cost.

Evaluation must reflect deployment. Rather than reporting only clean‑set accuracy, robust assessment exposes models to controlled uncertainty sources and measures both performance and calibration drift. For TSR, this includes corruption suites (noise, blur, weather, compression), crop perturbations to mimic detector imprecision, partial occlusions, and cross‑style or cross‑camera domain shifts. Because TSR classes can be visually similar (for example, speed‑limit signs), small degradations can disproportionately harm calibration even when accuracy declines modestly.

Finally, trustworthy TSR is a pipeline problem. Detector confidence determines which regions reach the classifier; classifier confidence summarizes class ambiguity in the crop; fusion and tracking can stabilize decisions across frames. Understanding how uncertainty propagates through these stages—and designing simple, fast interventions to keep confidences meaningful—is central to deploying TSR that remains reliable when conditions are less than ideal.

\subsection{Thesis area}

This thesis analyzes and improves the trustworthiness of traffic sign classifiers under uncertainties within a real‑time detection+classification pipeline. Using GTSRB as the primary benchmark, the work will create controlled uncertainty conditions—such as noise, blur, weather effects, compression, crop jitter to emulate detector errors, and partial occlusions at multiple severities—and quantify how each source affects both accuracy and calibration.

The study will train ResNet‑based classifiers in PyTorch as baselines and evaluate them with reliability diagrams, Expected Calibration Error, Brier score, and Negative Log‑Likelihood, alongside latency measurements to respect real‑time constraints. Lightweight improvements will be prioritized: post‑hoc calibration (temperature scaling and isotonic regression), single‑pass confidence proxies (maximum probability, entropy, energy), and selective prediction or abstention policies that maintain acceptable risk at target coverage.

Where compute permits, budget‑aware variants such as Monte‑Carlo dropout with a small number of passes or a compact ensemble may be included to compare trust‑per‑millisecond trade‑offs. The outcome will be a reproducible pipeline, a clear characterization of how different uncertainties degrade confidence quality, and practical recommendations—calibration settings, thresholds, and decision policies—for deploying TSR that remains reliable when inputs are imperfect, all while staying within real‑time limits.

\section{State of the Art}

Traffic Sign Recognition (TSR) in practice is a two-stage process: a real-time detector proposes candidate regions in the full image, and a classifier assigns a semantic label to each crop. Lightweight one‑stage detectors are favored on embedded hardware, while ResNet‑style backbones are common for classification \cite{stallkamp2011gtsrb,he2016deep,redmon2018yolov3}. On curated datasets such as GTSRB, clean, centered crops can yield high top‑1 accuracy \cite{stallkamp2011gtsrb}.

However, deployment settings introduce multiple sources of uncertainty: sensor noise, motion blur, adverse weather, compression artifacts, illumination changes, partial occlusions, crop jitter from imperfect detections, and domain shifts across cameras or regions. Under these conditions, confidence estimates often become unreliable: models may be confidently wrong, undermining downstream safety logic.

A large body of work addresses probability calibration. Post‑hoc methods such as temperature scaling and isotonic regression adjust the final scores using a held‑out calibration set and add essentially no runtime overhead \cite{guo2017calibration,niculescu2005predicting}. Training‑time techniques (e.g., label smoothing or mixup) can influence calibration as well, but their effects vary with dataset and model capacity. Diagnostics include reliability diagrams, Expected Calibration Error, Brier score, and Negative Log‑Likelihood, all of which quantify how predicted probabilities align with empirical correctness \cite{niculescu2005predicting,guo2017calibration}.

Uncertainty estimation complements calibration by indicating how much a prediction should be trusted. Single‑pass proxies such as maximum softmax probability and entropy are attractive for real‑time operation because they introduce negligible latency. Multi‑pass approaches—Monte‑Carlo dropout, stochastic weight sampling, or small deep ensembles—typically improve error–confidence alignment but increase inference time and memory \cite{gal2016dropout,lakshminarayanan2017simple}.

Decision‑making under uncertainty is often studied through selective prediction (classification with a reject option). By abstaining on low‑confidence inputs, systems can reduce error at the cost of coverage; risk–coverage curves summarize these trade‑offs \cite{geifman2019selectivenet}. Thresholds are typically set on a calibration set; lightweight conformal‑style thresholds can provide coverage guarantees under mild assumptions \cite{vovk2005algorithmic}. In TSR, selective strategies are especially useful when classes are visually similar (e.g., speed limits) and degradations remove fine details needed for disambiguation.

End‑to‑end pipeline effects remain underexplored. Many studies evaluate classification on perfect GTSRB crops, which hides failures introduced by detection, such as small or off‑center crops that erase critical digits. Detector confidence and non‑max suppression interact with classifier scores, so trust should be measured across the full pipeline and reported alongside latency. Robustness studies based on corruption benchmarks show that accuracy and calibration can degrade substantially under common corruptions, highlighting the need for stress testing beyond clean data \cite{hendrycks2019imagenetc}.

Reporting practices also vary widely. Some works present accuracy without calibration metrics; others report ECE but omit reliability plots, per‑severity breakdowns, or timing. Comparisons across uncertainty methods often do not normalize for computational cost, even though “trust per millisecond” is the key constraint in real‑time TSR.

\subsubsection*{Key limitations in existing solutions}

\begin{itemize}
  \item Calibration is frequently evaluated only on clean data; behavior under realistic uncertainties (noise, blur, weather, compression), detector crop jitter, and partial visibility is insufficiently characterized.
  \item Deployment‑friendly methods are undercompared; many results rely on ensembles or numerous stochastic passes that exceed real‑time budgets.
  \item Selective prediction and simple, guaranteed thresholding are rarely assessed specifically for TSR with clear coverage–risk and latency trade‑offs.
  \item End‑to‑end trust (detector + classifier) is seldom analyzed; assumptions of perfect crops mask failure modes introduced upstream.
  \item Reproducible stress testing on GTSRB lacks standardized multi‑severity perturbations with calibration and timing reported together.
\end{itemize}

This thesis addresses these gaps by building a reproducible, real‑time oriented evaluation that measures both performance and confidence quality under controlled uncertainties, compares lightweight improvements, and analyzes trust for the full detection+classification pipeline.

\section{Problem Statement}

\subsection{Thesis focus}

Despite strong accuracy on curated datasets, traffic sign recognition systems often produce overconfident errors once uncertainty is introduced by real-world conditions. In a detection+classification pipeline, uncertainty arises from sensor noise, motion blur, weather, compression, domain shift, imperfect detections (crop jitter, scale changes), small object size, and partial visibility. The central problem is that classifier confidences, which downstream components rely on, may no longer reflect true likelihoods of correctness, undermining safety and decision logic in real time.

This thesis addresses the problem by (i) quantifying how different uncertainty sources affect both performance and calibration in a realistic pipeline, and (ii) developing lightweight methods that restore trustworthy confidence without violating latency constraints.

The work will design and implement the following artifacts:

\begin{itemize}
  \item A reproducible PyTorch pipeline that integrates a lightweight detector with a ResNet-based classifier, including tools to inject controlled uncertainties (noise, blur, weather, compression, crop jitter, partial visibility) at multiple severities.
  \item A trust module that provides single-pass confidence measures (maximum probability, entropy, energy) and post-hoc calibration methods (temperature scaling, isotonic regression) with negligible overhead.
  \item Optional budget-aware uncertainty baselines (e.g., a small number of Monte-Carlo dropout passes or a compact ensemble) to study trust gains per millisecond of additional latency.
  \item A selective prediction component that sets abstention thresholds on a held-out calibration set to achieve target coverage with bounded risk; thresholds are chosen to be simple and deployable.
  \item An evaluation harness that computes accuracy, Negative Log-Likelihood, Brier score, Expected Calibration Error, reliability diagrams, risk–coverage curves, and end-to-end latency; results are reported per uncertainty type and severity.
\end{itemize}

Methodology and scope:

\begin{itemize}
  \item Dataset: GTSRB for classification; detector-generated crops and controlled crop perturbations emulate realistic upstream noise when ground-truth boxes are unavailable.
  \item Training: ResNet-18 (and optionally ResNet-34) baselines with standard augmentation and regularization; clean/occluded and corrupted evaluation splits.
  \item Uncertainty stress tests: additive noise, blur (motion/defocus), weather effects, compression, brightness/contrast changes, crop jitter, and partial visibility masks—each with multiple severities.
  \item Real-time constraint: profile detector+classifier inference; prioritize single-pass methods; treat multi-pass methods as optional comparisons with explicit latency budgets.
  \item Analysis: per-class and per-severity breakdowns, error–confidence alignment, and ablations on calibration set size, class imbalance, augmentation strength, and crop quality.
\end{itemize}

Expected outcome is a practical set of recommendations—calibration settings, confidence thresholds, and simple selection policies—that maintain trustworthy behavior under uncertainties while meeting real-time requirements.

\subsection{Research questions}

\textbf{Main question.} How trustworthy are traffic sign classifier confidences in a real-time detection+classification pipeline under common uncertainties, and which lightweight methods best restore calibration and decision quality without exceeding latency budgets?

\begin{enumerate}
  \item \textbf{Explanation:} Why do specific uncertainty sources (e.g., noise, blur, crop jitter, partial visibility) degrade calibration more than others in TSR?
  \item \textbf{Description:} What do the accuracy, calibration (ECE, Brier, NLL), and reliability diagrams look like across uncertainty types and severities for ResNet-based baselines?
  \item \textbf{Assessment:} How do post-hoc methods (temperature scaling, isotonic regression) and single-pass confidence proxies (max probability, entropy, energy) compare in risk–coverage trade-offs per millisecond of overhead?
  \item \textbf{Assessment:} Can simple abstention thresholds, fitted on a small calibration set, maintain target coverage with acceptable risk across uncertainty shifts?
  \item \textbf{Prediction/Assessment:} Given a fixed latency budget, what incremental benefit is obtained from small-budget stochastic methods (e.g., few-pass MC dropout or a compact ensemble), and is the gain justified relative to their cost?
  \item \textbf{Description/Assessment:} How sensitive are trust metrics to detector imperfections (crop tightness, off-center crops), and can post-hoc calibration compensate for such upstream errors?
\end{enumerate}

These questions are designed to be answerable within the thesis by the proposed pipeline, metrics, and experiments, and together they cover the core research problem of trustworthy TSR under uncertainties.

\section{Approach}

The thesis will begin by establishing a simple, reproducible detection+classification pipeline and a controlled uncertainty benchmark, then layering in lightweight calibration and decision policies. The early focus is on single‑pass methods that add negligible latency; multi‑pass methods are treated as optional comparisons to quantify trust‑per‑millisecond trade‑offs.

The pipeline will be implemented in PyTorch using GTSRB for classification. A fixed data split will reserve a small calibration set for post‑hoc score adjustment and threshold selection. To reflect deployment, uncertainties will be introduced through a modular “stress suite” with multiple severities: additive noise, motion/defocus blur, brightness/contrast changes, weather overlays, JPEG compression, and crop jitter to emulate detector imprecision. Partial visibility masks can be included as one uncertainty source but the analysis remains general to “under uncertainties.”

Promising ideas to explore include:
\begin{itemize}
  \item \textbf{Single‑pass trust signals:} maximum probability, entropy, and energy‑based scores as fast proxies.
  \item \textbf{Post‑hoc calibration:} temperature scaling and isotonic regression fitted on the calibration set; per‑class versus global variants compared for stability.
  \item \textbf{Selective prediction:} simple abstention thresholds tuned to target coverage, plus a lightweight conformal-style threshold for coverage guarantees.
  \item \textbf{Budgeted baselines (optional):} few‑pass Monte‑Carlo dropout or a compact ensemble to bound potential gains relative to latency.
  \item \textbf{Pipeline realism:} integrate a lightweight pre‑trained detector to generate crops; when unavailable, simulate detection noise via controlled crop jitter to study its impact on calibration.
\end{itemize}

Evaluation will report accuracy, Negative Log‑Likelihood, Brier score, Expected Calibration Error, reliability diagrams, risk–coverage curves, and end‑to‑end latency. Analyses will include per‑severity breakdowns, error–confidence alignment, and ablations on calibration set size, augmentation strength, and crop quality. All experiments will use fixed seeds and configuration files for reproducibility, with automated plotting and table generation.

\textbf{First steps after the proposal (weeks 1–3):}
\begin{enumerate}
  \item Set up repository, data loaders, and fixed train/validation/test splits; carve out a calibration set.
  \item Train a baseline ResNet‑18 on clean GTSRB crops; implement metrics and reliability plots.
  \item Implement the uncertainty stress suite (noise, blur, weather, compression, crop jitter) with severity levels and re‑run baseline evaluation.
  \item Add temperature scaling and isotonic regression; compare single‑pass confidence measures and generate initial risk–coverage curves.
\end{enumerate}

Subsequent weeks will integrate the detector front‑end, perform latency profiling, and, if time allows, add budgeted multi‑pass baselines to complete the trust‑per‑millisecond comparison.

\section{Planning}

\subsection{Own Background}

I have practical experience with Python, PyTorch, and probability/statistics, as well as foundational knowledge of linear algebra, calculus, and machine learning. I am comfortable with CNN training, data preprocessing, and basic experiment tracking and plotting in Jupyter/Matplotlib.

\subsection{Required Resources}

\textbf{Hardware}

\begin{itemize}
  \item 1 GPU with $\geq$ 8--12\,GB VRAM (e.g., RTX 3060/3070 or university server GPU) for training and evaluation.
  \item Multi-core CPU, 16--32\,GB RAM, and $\sim$20\,GB storage for datasets, checkpoints, and logs.
  \item Reliable backup/storage and internet access for package and model downloads.
\end{itemize}

\textbf{Software}

\begin{itemize}
  \item Linux, Python 3.10+, PyTorch and TorchVision; optional TensorFlow for cross-checks.
  \item Data/augmentation: Albumentations or \texttt{imgaug}; metrics: scikit-learn; plotting: Matplotlib/Seaborn.
  \item Experiment tracking: TensorBoard or Weights \& Biases; profiling: \texttt{torch.profiler}.
  \item Pre-trained lightweight detector (public weights) for sign region proposals.
\end{itemize}

\textbf{Data and assets}

\begin{itemize}
  \item GTSRB (classification) with a held-out calibration split.
  \item Synthetic uncertainty components (noise, blur, weather overlays, compression) and crop-jitter tooling.
  \item Optional: public traffic-sign detector model and any auxiliary scripts for crop extraction.
\end{itemize}

\subsection{Work packages}

Dated timeline for 13 Apr 2026 -- 15 Aug 2026 (single-student, \(\approx\)13 h/week). Key milestones are marked in brackets.

\begin{enumerate}[label=\textbf{W\arabic*}, leftmargin=*]
  \item 13--19 Apr: Kickoff — finalize scope/RQs, set up repo/environment, download GTSRB, initial reading.
  \item 20--26 Apr: Project definition — confirm milestones and latency targets; refine proposal with supervisor; focused literature map. \textbf{[M1: Proposal confirmed]}
  \item 27 Apr--03 May: Research \& planning — choose toolchain (PyTorch), configs, CI; freeze splits incl.\ a held‑out calibration set.
  \item 04--10 May: Data handling — clean/verify GTSRB, EDA, class balance, preprocessing; document train/val/test/calibration protocol.
  \item 11--17 May: Baseline modeling — train ResNet‑18 on clean crops; implement accuracy, NLL, Brier, ECE, reliability diagrams. \textbf{[M2: Baseline ready]}
  \item 18--24 May: Model development — regularization/augmentation tuning; stabilize baseline; seed control and logging.
  \item 25--31 May: Uncertainty stress suite v1 — add noise/blur/weather/compression (multi‑severity); evaluate accuracy \& calibration drift.
  \item 01--07 Jun: Mid‑term review — present baseline + stress results; adjust plan/ablations with supervisor. \textbf{[M3: Mid‑term review]}
  \item 08--14 Jun: Detection realism — implement crop jitter; integrate a lightweight pre‑trained detector to generate crops; compare to perfect crops.
  \item 15--21 Jun: System integration — end‑to‑end pipeline; latency profiling; optimizations to meet real‑time constraints. \textbf{[M4: Real‑time E2E achieved]}
  \item 22--28 Jun: Calibration \& trust — add temperature scaling \& isotonic regression; compare single‑pass proxies (max‑prob, entropy, energy); start selective prediction.
  \item 29 Jun--05 Jul: Ablations \& advanced tests — per‑class analyses; sensitivity to crop tightness/off‑center; optional few‑pass MC dropout or small ensemble if time allows. \textbf{[M5: Results freeze]}
  \item 06--12 Jul: Presentation preparation — finalize plots/tables; compile recommendations (thresholds, policies); draft slides and demo.
  \item 13--19 Jul: Writing block — methods, experiments, results; supervisor feedback round.
  \item 20--26 Jul: Wrap‑up draft — discussion, limitations, future work; polish figures; prepare artifact package (code/configs).
  \item 27 Jul--02 Aug: Buffer \& revisions — address feedback; dry‑run talk; finalize timings and demo assets.
  \item 03--09 Aug: Final presentation/defense window — deliver talk and demo; incorporate committee feedback. \textbf{[M6: Defense]}
  \item 10--15 Aug: Final wrap‑up — camera‑ready thesis; artifact release; checklist and submission. \textbf{[M7: Final submission]}
\end{enumerate}

\subsection{Contingency plan}

\begin{itemize}
  \item \textbf{Risk:} Real-time budget is exceeded by multi-pass methods. \textbf{Mitigation:} Prioritize single-pass calibration and confidence proxies; keep MC dropout/ensembles optional and report trust gains normalized by latency.
  \item \textbf{Risk:} Calibration fails to generalize across uncertainty types. \textbf{Mitigation:} Compare global vs. per-class temperature, isotonic regression, and threshold recalibration per uncertainty family; use simple conformal-style thresholds to control coverage.
  \item \textbf{Risk:} Compute or time constraints limit experimentation. \textbf{Mitigation:} Freeze scope at single-pass methods and the most impactful uncertainties; reduce grid sizes; reuse pre-trained backbones.
  \item \textbf{Risk:} Results are hard to reproduce. \textbf{Mitigation:} Fixed seeds, config files, environment locking, and automated report generation; publish scripts and documentation.
\end{itemize}

\begin{thebibliography}{10}

\bibitem{stallkamp2011gtsrb}
J.~Stallkamp, M.~Schlipsing, J.~Salmen, and C.~Igel,
\newblock ``The German Traffic Sign Recognition Benchmark: A multi-class classification competition,''
\newblock In Proc. IJCNN, 2011.

\bibitem{he2016deep}
K.~He, X.~Zhang, S.~Ren, and J.~Sun,
\newblock ``Deep Residual Learning for Image Recognition,''
\newblock In Proc. CVPR, 2016.

\bibitem{redmon2018yolov3}
J.~Redmon and A.~Farhadi,
\newblock ``YOLOv3: An Incremental Improvement,''
\newblock arXiv:1804.02767, 2018.

\bibitem{guo2017calibration}
C.~Guo, G.~Pleiss, Y.~Sun, and K.~Q. Weinberger,
\newblock ``On Calibration of Modern Neural Networks,''
\newblock In Proc. ICML, 2017.

\bibitem{niculescu2005predicting}
A.~Niculescu-Mizil and R.~Caruana,
\newblock ``Predicting Good Probabilities with Supervised Learning,''
\newblock In Proc. ICML, 2005.

\bibitem{gal2016dropout}
Y.~Gal and Z.~Ghahramani,
\newblock ``Dropout as a Bayesian Approximation: Representing Model Uncertainty in Deep Learning,''
\newblock In Proc. ICML, 2016.

\bibitem{lakshminarayanan2017simple}
B.~Lakshminarayanan, A.~Pritzel, and C.~Blundell,
\newblock ``Simple and Scalable Predictive Uncertainty Estimation using Deep Ensembles,''
\newblock In Proc. NeurIPS, 2017.

\bibitem{geifman2019selectivenet}
Y.~Geifman and R.~El-Yaniv,
\newblock ``SelectiveNet: A Deep Neural Network with an Integrated Reject Option,''
\newblock In Proc. ICML, 2019.

\bibitem{vovk2005algorithmic}
V.~Vovk, A.~Gammerman, and G.~Shafer,
\newblock \emph{Algorithmic Learning in a Random World},
\newblock Springer, 2005.

\bibitem{hendrycks2019imagenetc}
D.~Hendrycks and T.~Dietterich,
\newblock ``Benchmarking Neural Network Robustness to Common Corruptions and Perturbations,''
\newblock In Proc. ICLR, 2019.

\end{thebibliography}

\end{document}